% ------------------------------------------------------------
% Page 87
% ------------------------------------------------------------

\section*{En remontant la vallée de la Brèze}

Nous remontons doucement la route qui suit la rivière et après quelques grands gouffres où il
fait bon se baigner par grosse chaleur, s’ouvre sur la gauche le petit vallon de Cabanals.

\textbf{Les jardins :} plusieurs familles de Meyrueis y possèdent des parcelles exiguës, bien abritées
du vent du nord et les bancels inférieurs « étaient à l’arrosage », car de petites chaussées
construites dans le ruisseau alimentaient des « béals ».

\textbf{Les châtaigniers :} dans cette vallée schisteuse, prospéraient les premiers châtaigniers,
ressource inestimable pour faire la soupe ou accompagner, les jours de fête, un morceau de
porc.

\textbf{Les ardoises} qu’on pouvait extraire étaient de bonne qualité et couvraient bien des toits de
Meyrueis et des environs puisque, on en « exportait » sur le causse et même dans les
départements voisins : Aveyron ou Gard qui ne sont pas loin. En 1880, Henri et Maurice
Séquier, couvreurs de Meyrueis employaient trois compagnons à Cabanals et « sortaient aussi
des ardoises dans les chênes entre Raffègues et Alauze ».

\textbf{L’oxyde de manganèse} fut exploité entre 1830 et 1840 environ pour les hauts fourneaux
d’Alès.

\vspace{0.2cm}

Au sortir du vallon apparaît \textbf{Alauze} et son pigeonnier. Au 13° siècle la ferme était, tout
comme Campis ou Sérigas propriété du chapitre de l’Eglise Cathédrale de Montpellier, prieur
d’Ayres qui l’aliéna en 1578. Plus tard, bien plus tard, en 1942, c’est au bord de la Brèze que
fut construite l’usine de carburants forestiers, l’usine d’Alauze, qui fonctionna quelques
années.

\vspace{0.2cm}

Sur la droite de l’autre côté de la rivière, voilà les bâtiments de \textbf{Raffègues.}

Dans les années 1790 un certain Antoine Maillé était propriétaire ou fermier. Plainte est
déposée contre lui en 1799 « contre Antoine Maillé du domaine de Raffègues qui autorisa
dans un bâtiment de son domaine un rassemblement de fanatiques pour l’exercice du culte
(11 Thermidor-17 Fructidor an VII (Août --Sept 1799)). ».*

S’agit-il d’une conséquence de la Terreur Blanche ? Sans doute.

\vspace{0.2cm}

Mais un peu plus tard à Raffègues toujours : « On assiste, écrit Ph. Chambon, à d’authentiques
actes de solidarité comme, par exemple, la cachette qu’Esaïe Combes, fermier de Raffègues et
notable protestant, offre à M. Trophime Bragouse de Saint Sauveur, ancien grand vicaire de
l’archevêque de Narbonne et prêtre insermenté.

Un autre protestant en vue à l’époque Antoine Avesque qui s’était engagé en 1792 avec 88
autres jeunes de Meyrueis pour défendre la République aux frontières, permettra au dit grand
vicaire de célébrer la messe clandestinement, dans une grange de son mas de Ribevenès ».

\vspace{0.3cm}

\begin{itemize}
    \item Archives de la Lozère : C 253
\end{itemize}
